% GENERATED FILE. Edit canonical lessons, then run npm run generate:books.
% canonical-source-hash: fnv1a64:ded797c79a7adb90
% canonical-lessons: TE-C14-rutuvulu, TE-S118-letter-ga

\chapter{Seasons}
\label{ch:te-seasons}
\hlchaptermodality{14}

\begin{chapteropening}
\textbf{By the end of this chapter:} \emph{I can name the four seasons in Telugu, including the rains that anchor the year.}

Say \te{వేసవి} (cousin of Kannada's bēsige), \te{వానాకాలం} and \te{శీతాకాలం} from native roots, then the doubly-Sanskrit \te{వసంత} \te{ఋతువు}.
\end{chapteropening}

\section[\texorpdfstring{\te{వసంత} \te{ఋతువు} \te{వేసవి} … (vasanta ṛtuvu …)}{వసంత ఋతువు వేసవి … (vasanta ṛtuvu …)}]{\te{వసంత} \te{ఋతువు}, \te{వేసవి}, \te{వానాకాలం}, \te{శీతాకాలం} — the pattern completes itself}

\label{lesson:TE-C14-rutuvulu}



\begin{quote}
Telugu closes out the Dravidian season words the same way Kannada did: native words for summer/rain/cold, and a spring-word borrowed twice over from Sanskrit.
\end{quote}

\subsection*{You'll want to know: The words}
\noindent
\begin{tabularx}{\linewidth}{@{}>{\raggedright\arraybackslash}X>{\raggedright\arraybackslash}X>{\raggedright\arraybackslash}X@{}}
\toprule
\textbf{Telugu} & \textbf{meaning} & \textbf{source} \\
\midrule
\textbf{\te{వసంత} \te{ఋతువు}} (\emph{vasanta ṛtuvu}) & \textquotedblleft{}spring\textquotedblright{} & both \emph{vasanta} and \textbf{\te{ఋతువు}} \emph{ṛtuvu} (\textquotedblleft{}season\textquotedblright{}) are Sanskrit, like Kannada's pair \\
\textbf{\te{వేసవి}} (\emph{vēsavi}) & \textquotedblleft{}summer\textquotedblright{} & native Dravidian — cousin of Kannada's \emph{bēsige} \\
\textbf{\te{వానాకాలం}} (\emph{vānākālam}) & \textquotedblleft{}rainy season\textquotedblright{} & \textbf{\te{వాన}} \emph{vāna} (\textquotedblleft{}rain,\textquotedblright{} native) + \emph{kālam} (\textquotedblleft{}time,\textquotedblright{} itself Sanskrit-derived but long assimilated) \\
\textbf{\te{శీతాకాలం}} (\emph{śītākālam}) & \textquotedblleft{}winter/cold season\textquotedblright{} & \emph{śīta} (\textquotedblleft{}cold\textquotedblright{}) leans Sanskrit-flavored here, unlike Tamil/Kannada's fully native cold-words \\
\bottomrule
\end{tabularx}

\begin{culture}
Telugu's \emph{vēsavi} and Kannada's \emph{bēsige} are close enough in shape to be genuine cousins — the same native Dravidian summer-word wearing two languages' sound changes. And once again, it's \textbf{\te{వానాకాలం}} (\emph{vānākālam}, the rains) that actually organizes Andhra/Telangana's agricultural year, not a clean four-way Western split — the same honest point its neighbours make.
\end{culture}

\subsection*{Your turn}
Say these aloud:

\begin{itemize}
\raggedright
  \item \textquotedblleft{}vēsavi\textquotedblright{} — summer, cousin of Kannada's bēsige
  \item \textquotedblleft{}vānākālam\textquotedblright{} — the rains, the real central season
  \item \textquotedblleft{}śītākālam\textquotedblright{} — the cold season
  \item the double Sanskrit loan — \textquotedblleft{}vasanta ṛtuvu,\textquotedblright{} like Kannada
\end{itemize}

\begin{tcolorbox}[breakable,colback=teal!4,colframe=teal!35!black,title={Before you move on}]
What's the relationship between Telugu's \emph{vēsavi} and Kannada's \emph{bēsige}? (\textbf{Cousins} — the same native Dravidian summer-word.) Which Dravidian languages borrow BOTH halves of their spring-word from Sanskrit (the season-name AND the word for \textquotedblleft{}season\textquotedblright{} itself)? (\textbf{Kannada and Telugu} — \emph{vasanta ṛtu(vu)} — unlike Tamil/Malayalam's single \emph{vasantham} loan plus the long-assimilated \emph{kaalam}.)
\end{tcolorbox}
\section[\texorpdfstring{\te{గ} (ga)}{గ (ga)}]{\te{గ} — one character, met inside words you already say}

\label{lesson:TE-S118-letter-ga}



\begin{quote}
Before the new one: \te{వ} — what does it do? And one from further back: \te{ల}?

One character this time — and you have been saying it for pages without knowing which mark on the page it was.
\end{quote}

\subsection*{Script you'll notice: \te{గ}}
\textbf{\te{గ}} — \emph{ga}.

It is a \textbf{consonant}, and in this script a consonant is never bare: it comes with an \emph{a} already in it. So it is not \emph{g}, it is \textbf{ga}.

You already say these, and every one of them has \te{గ} somewhere inside it:

\begin{itemize}
\raggedright
  \item \textbf{\te{బాగా}} \emph{bāgā} — well, fine — and the reply "\te{నేను} \te{బాగున్నాను}
  \item \textbf{\te{నేను} \te{తెలుగు} \te{మాట్లాడతాను}} — I speak Telugu
  \item \textbf{\te{నాకు} \te{తెలుగు} \te{వచ్చు}} \emph{nāku telugu vaccu} — 'I know Telugu'
  \item \textbf{\te{ఒకటి} \te{రెండు} \te{మూడు} \te{నాలుగు} \te{ఐదు}} \emph{okaṭi reṇḍu mūḍu nālugu aidu} — one to five
\end{itemize}

\subsection*{Writing: \te{గ} — follow the numbered strip}
\hlblockfigure{\detokenize{figures/TE-S118-letter-ga-filmstrip.pdf}}{How \te{గ} is written, stroke by stroke}

Follow the numbered strip: start where movement 1 begins, and let each panel show you the next movement before your pen makes it. Then write \te{గ} yourself — slowly, and larger than it is printed.

\begin{quote}
The strip shows \textbf{where to start the character and which way to travel}, in one attested order. That is taught with real variation from school to school, so the strip names its source beneath it: treat it as a sound way in, not the only one.
\end{quote}

\subsection*{Your turn}
\begin{itemize}
\raggedright
  \item \emph{Look:} at these words, and find \te{గ} in the ones that have it
\end{itemize}

\begin{quote}
\te{బాగా}  ·  \te{నమస్కారం}
\end{quote}

\begin{itemize}
\raggedright
  \item \emph{Trace it:} \te{గ} three times, saying \emph{ga} as you finish each one
  \item \emph{Look:} back at any page of this chapter and find \te{గ} once more
\end{itemize}

\begin{tcolorbox}[breakable,colback=teal!4,colframe=teal!35!black,title={Before you move on}]
Which character is this — \te{గ}? What sound does it carry? (\emph{\textbf{ga}}.) Name one word you already say that contains it.

One more, from much earlier: \te{స} — what does it do?
\end{tcolorbox}
